% s-201-p201-m002; original PDF 201
\Aphorism{s-201-p201-m002}{286}
Es ist ein Irrthum, dass der Dichter sich völlig und explicit mittheilen solle oder auch nur dürfe, mit allen Nüancen, Feinheiten und Nebenbeziehungen. Er hat nur für die Klangfarbe zu sorgen, für das Mitschwingen der Obertöne zu seinem Grundton; aber dies Ziel würde er gründlich verfehlen, wenn er alle Obertöne wirklich anschlagen wollte. Er muss sein Instrument kennen und seinen Resonanzboden, dann weiss er, \Sperr{welchen }Ton er greifen muss, um \Sperr{welchen anderen} zu Gehör zu bringen! Allerdings kann man, was den Resonanzboden „Publicum“ anbelangt, wunderliche Dinge erleben; es liegt eine Art Raffinement darin, wie der Leser jede nicht ganz plane Deutlichkeit des Autors missversteht. Wenn sich heutzutage keine Prosa mehr findet, die in \Sperr{anmuthiger} Weise (von der unterschwürig-geheimnissthuerischen Drei-Punkte-Poesie rede ich nicht) mit Nebensinn, Anspielung, doppeltem Boden zu wirken weiss, so ist das heutige Publicum dafür verantwortlich, das eilfertige telegraphirende und telephonirende \Foreign{„time is money}“-Publicum, dem man auch in der Kunst nur die drastische Eindeutigkeit eines Cursberichts bieten darf. Die Zeit der zierlich gesponnenen Conversation, des Federballspiels zwischen Frage und Antwort ist vorüber; aber sie kommt wieder, sie kündigt sich bereits in einigen kleinen und exquisiten Arbeiten jüngerer Franzosen an, auch in Ibsens letzten Dramen, wo das Zusammenklingen von Sinn und Nebensinn, die „Transparenz“ der Sprache, das Mittönen eines unhörbaren Dialogs das Entzücken feiner Ohren ist.
\par

